\section{Conclusions}
\label{sec:conclusions}

\status{Depends on the sections above. What can be written now is listed; the
rest is marked.}

\missing{conclusions --- to be written once the matrix is closed}

\subsection*{What is established}

\begin{itemize}
\item Both methods reproduce the ellipsoidal regime of the Newtonian reference
case and the shape of its published rise-velocity curve, and they \emph{converge
toward each other} under refinement --- from 1.94\% to 0.60\% in centre of mass
and from 4.3\% to 1.1\% in terminal rise velocity --- approaching from opposite
sides.
\item At the reference resolution \textbf{VOF} conserves mass about five times
better ($-0.170\%$ against $-0.851\%$), and front-tracking's drift \emph{grows}
with refinement. The coarse-mesh figures reverse this ranking and are partly
cancellation; they should not be quoted.
\item The two methods cost the same to within 5\% per cell-step, with the sign
of the difference depending on resolution.
\item The verification apparatus separates exact from convergent oracles, and
the exact ones --- surface closure, genus, the plane/slab/tetrahedron fraction
--- pass at machine precision.
\item The viscoelastic matrix reproduces the reference results. All twelve
ordering claims of the reference thesis hold, the quoted values agree to within
7\% (most within 3\%), the trend against a Newtonian reference in the same
geometry holds in all four cases, and the specific velocity signature the
reference describes --- a sharp drop after the maximum, then recovery ---
appears in both methods.
\item \textbf{VOF sits above the published values on all nine comparisons and
front-tracking at or below on seven of nine}, with mean absolute deviations of
$3.6\%$ and $1.2\%$. The asymmetry is the opposite of what the provenance
suggests, since the reference was computed with VOF.
\item Oldroyd-B did \emph{not} run away on the time scale of the rise: all
twelve runs completed 2001 steps with zero non-finite values.
\end{itemize}

\subsection*{What is not}

\begin{itemize}
\item Adaptive refinement in production. The dynamic remeshing loop works and
is verified, but no case has been run on an adapted mesh to completion.
\item Why front-tracking's mass drift grows with refinement. The triangulation
counts and the surgery counts are nearly identical at both resolutions, so
neither explains it.
\item Which of the two readings explains the sign asymmetry: a difference
between our volume-of-fluid and the reference's, or an upward bias in the
marching-cubes extraction from which the shape quantities are read.
\item Case~2, where topological change is expected.
\item Any statement about more than one process.
\end{itemize}
